\documentclass[11pt,letterpaper]{article}
\usepackage[margin=1in]{geometry}
\usepackage[T1]{fontenc}
\usepackage{lmodern}
\usepackage{microtype}
\usepackage{amsmath,amssymb}
\usepackage[numbers,sort&compress]{natbib}
\usepackage{booktabs}
\usepackage[skip=6pt]{caption}
\usepackage{enumitem}
\usepackage[compact]{titlesec}
\usepackage[hidelinks]{hyperref}
\hypersetup{pdftitle={Does an Edge-of-Stability Mechanism Explain E3M4 versus E2M5?}}
\setlength{\parindent}{0pt}
\setlength{\parskip}{0.35em}
\setlength{\emergencystretch}{2em}
\clubpenalty=10000
\widowpenalty=10000
\setlist{nosep,leftmargin=1.2em}

\title{\vspace{-1.6em}Does an Edge-of-Stability Mechanism Explain E3M4 versus E2M5?\\[2pt]
\large Literature review and implications (9.520, Fall 2026)}
\author{Juntang Wang \and V\'ictor Conchello}
\date{October 2, 2026}

\begin{document}
\maketitle
\vspace{-2em}

\section{The question and why it is hard}

An 8-bit E$x$M$y$ format spends $x$ bits on the exponent and $y$ on the mantissa, so E3M4 has more range than E2M5 (about $2^{11}$ against $2^{8}$ under the OFP8 convention \citep{Micikevicius2022}) but coarser precision (rounding error $2^{-5}$ against $2^{-6}$). As the shared scales of FP4 blocks, E3M4 gives a lower LLM training loss \citep{Chmiel2025}, and nobody has explained why. We ask whether the edge of stability (EoS) does: training sits at the stability limit of the sharpest curvature direction, so quantization error along it could matter far more than its average size suggests. This is hard to test. Range and precision change together, the exponent bias alone can reverse a ranking \citep{Kuzmin2022}, and no study has measured sharpness while varying the format.

\section{E3M4 versus E2M5}

\subsection{How the literature distinguishes the two formats}

\paragraph{\citet{Chmiel2025}.}
\emph{E3M4 works better than E2M5 as the scale format for FP4 training.} This study trains LLMs with FP4 (E2M1) weights, activations and gradients, where each block of 16 values shares one 8-bit scale that restores its real magnitude in higher-precision computation. Across seven simulated scale formats on a 350M Llama-style model, E3M4 and E4M3 reach the lowest loss; E2M5, E5M2 and E8M0 end slightly worse, E6M1 clearly worse, and E1M6 diverges. This challenges the Microscaling (MX) standard \citep{Rouhani2023}, whose E8M0 scale has no mantissa bits. A scale error rescales its whole block, so errors are correlated rather than independent. The ranking is purely empirical: the authors do not explain it.

\paragraph{\citet{Kuzmin2022}.}
\emph{More exponent bits pay off when tensors have outliers; more mantissa bits pay off otherwise.} This paper derives the expected quantization error of FP8 formats as a function of the tensor distribution: for Gaussian tensors E2M5 is optimal, and the optimum moves to more exponent bits as outliers grow. Pretrained FP32 transformers have heavy-tailed activations, so their post-training quantization favours more exponent bits. With a bias set per channel, E3M4 beats E2M5 on ViT (77.7 vs 77.3) and on BERT (82.6 vs 80.3); with a single bias for the whole network, E2M5 falls to 51.1 on ViT against 77.0 for E3M4. For CNNs the ranking depends on the bias. Quantization-aware training narrows the gaps between formats to about one accuracy point or less, but does not close them. This is a first hint that a format gap can partly reflect a mismatch between the training and deployment arithmetic.

\paragraph{Role of the format, by tensor.}
Early FP8 training work already chose the format according to the tensor: precision for weights and activations, range for gradients. \citet{Sun2019} used a 1-4-3 format in the forward pass and 1-5-2 in the backward pass. \citet{Noune2022} found that the exponent bias must be tuned, with more mantissa bits for weights and activations and more range for gradients. \citet{Micikevicius2022} proposed E4M3 and E5M2 with per-tensor scaling, later standardized by OCP as OFP8.

\paragraph{Distribution and outliers.}
Outlier features that grow with model scale \citep{Dettmers2022} widen the spread of block maxima, which favours range (E3M4). For well-behaved inference tensors, \citet{vanBaalen2023} find INT8 often preferable to FP8, and E2M5 is the more integer-like of the two formats, with fewer and longer binades. Format effects can also appear late. \citet{Fishman2025} trace FP8 divergences in trillion-token runs to outlier amplification in SwiGLU, which short sweeps such as Chmiel et al.'s would miss.

\paragraph{Rounding, accumulation and saturation.}
Round-to-nearest loses updates smaller than half a grid step, while stochastic rounding is unbiased; with 16-bit fixed point, only stochastic rounding matches FP32 training \citep{Gupta2015}. Neither is better everywhere: testing each tensor separately, \citet{Chmiel2025} find that stochastic rounding lowers the loss for neural gradients but raises it in the forward pass. Saturation is a separate failure mode. \citet{Su2025} trace loss spikes in microscaling training to values clamped at the top of the scaled range, and show that at 8 bits E4M3 beats E5M2. In both settings, more mantissa bits help as long as range suffices.

\section{The edge of stability}

\paragraph{\citet{Cohen2021}.}
\emph{Gradient descent drives sharpness up to the stability limit $2/\eta$ and then keeps it there while the loss still falls.} For a quadratic with curvature $\lambda$, a gradient step multiplies the error by $1-\eta\lambda$, so it converges only if the sharpness $\lambda_{\max}(\nabla^2L)<2/\eta$. Cohen et al.\ show that full-batch gradient descent violates the assumption behind classical analyses. Sharpness first rises (\emph{progressive sharpening}), then settles at or just above $2/\eta$. The loss oscillates along the top Hessian eigenvector $q_1$ but still decreases over long horizons. The next top eigenvalues follow, rising to $2/\eta$ one after another and staying there. With heavy-ball momentum the threshold becomes $(2+2\beta)/\eta$. Under SGD, sharpness does not settle at the threshold.

\paragraph{\citet{Damian2023}.}
\emph{Oscillations at the edge of stability reduce sharpness automatically, so training stays stable.} If the iterate moves a distance $x$ along $q_1$, the cubic Taylor term adds about $\tfrac{x^2}{2}\nabla^3L[q_1,q_1]=\tfrac{x^2}{2}\nabla\lambda_{\max}$ to the gradient, so gradient descent moves along $-\nabla\lambda_{\max}$. Oscillations along $q_1$ therefore reduce sharpness until the iterates are stable again. In aggregate, gradient descent follows projected gradient descent on $\{\theta:\lambda_{\max}(\theta)\le2/\eta,\ \nabla L\cdot q_1=0\}$. With several eigenvalues at the edge, training becomes more chaotic. The authors note that the dynamics are highly sensitive to small perturbations and switch to 64-bit gradients once training reaches instability, to avoid propagating floating-point error. That makes numerical precision a direct concern for any EoS-based explanation.

\paragraph{\citet{Cohen2022}.}
\emph{Adam also trains at the edge of stability, measured on the preconditioned sharpness.} For Adam with preconditioner $P=\mathrm{diag}(\sqrt{\nu})+\epsilon I$, where $\nu$ is the running average of squared gradients, the relevant quantity is the preconditioned sharpness $\lambda_{\max}(P^{-1}H)$. With $P$ held fixed, Adam on a quadratic diverges once it exceeds $2(1+\beta_1)/(\eta(1-\beta_1))=38/\eta$ for $\beta_1=0.9$, the analogue of $2/\eta$. In full-batch experiments, the preconditioned sharpness rises to $38/\eta$ and hovers there while the loss keeps falling, just as sharpness does at $2/\eta$ under gradient descent. AdamW is the standard optimizer for LLM training, Chmiel et al.'s included, so this paper is what makes the edge of stability relevant to modern LLM training.

\paragraph{Earlier evidence and linear stability.}
\citet{Wu2018} showed that gradient descent can only remain at minima with sharpness $\le2/\eta$, and SGD faces an additional constraint from gradient-noise non-uniformity. \citet{Jastrzebski2020} found that, early in training, SGD reaches a point where its steps become too large for the sharpest direction, which they call the \emph{break-even point}.

\paragraph{Theory of why training continues.}
\citet{Arora2022} prove that, for normalized gradient descent and for gradient descent run on the transformed loss $\sqrt{L-\min L}$ instead of the loss $L$ itself, the iterates enter the edge of stability and follow a sharpness-reducing flow on the manifold of minimizers. \citet{Ahn2022} characterize unstable convergence and the conditions that make it possible. \citet{Kreisler2023} prove that, in scalar networks, gradient descent monotonically decreases the sharpness of the gradient-flow solution it would reach.

\subsection{Has the edge of stability been used to compare number formats?}

We found no study that measures sharpness while varying number format. The closest is \citet{Su2025}. They model the low-precision gradient as $(1+\zeta_t)\bar g_t$ and derive a heuristic stability condition $|1-\eta\lambda_{\max}|+\eta\|\zeta_t\|\lambda_{\max}\lesssim1$. When $\eta\lambda_{\max}>1$, this shrinks the stable step size to $2/((1+\|\zeta_t\|)\lambda_{\max})$. They estimate a lower bound on $\|\zeta_t\|$ from paired FP32 and MXFP8 runs and find that the loss diverges once it reaches about 2. They frame this as a smaller EoS region, but never measure sharpness, so they do not check it against \citet{Cohen2021}. \citet{Wortsman2023}, studying int8 training of CLIP models, trace loss spikes to AdamW second-moment estimates that underestimate the squared gradients (mainly in the patch-embedding layer), and note that these spikes come with large activations and gradients that strain low-precision formats. That is a preconditioner-based account, not a format-specific curvature measurement. An EoS explanation of the format gap would require E2M5 to perturb training along $q_1$, or to bias $P$, more than E3M4 does. No study has tested this.

\subsection{A complementary mechanism: training--inference mismatch}

\paragraph{\citet{Qi2025}.}
\emph{Low precision destabilizes RL fine-tuning through a mismatch between the training and inference engines, and FP16 fixes it.} In RL fine-tuning of LLMs, the rollout engine (e.g.\ vLLM) and the training engine compute different token probabilities for the same weights, which biases the policy ratio and can collapse training. Algorithmic corrections only partially help, whereas switching both engines from BF16 to FP16 (7 versus 10 mantissa bits) largely removes the problem across algorithms and frameworks. \citet{Yao2025} document the same gap and correct it with truncated importance sampling, and \citet{He2025} show that the mismatch can arise even at fixed precision, from inference kernels that are not batch-invariant.

\section{Three key papers}

\paragraph{\citet{Cohen2021}} (ICLR 2021; about 500 citations on Semantic Scholar) founded the edge-of-stability literature. It supplies the mechanism we test and the quantity to measure: if stability explains the format gap, E3M4 and E2M5 must settle at different sharpness for the same learning rate.

\paragraph{\citet{Chmiel2025}} (NeurIPS 2025; about 45 citations in its first year) is the main training-time comparison of E3M4 and E2M5, so it defines the result we must explain. Its limits, one seed per format, one model size and no explanation, are what our experiments must address.

\paragraph{\citet{Kuzmin2022}} (NeurIPS 2022; about 150 citations) is the main analytic account of how the split between exponent and mantissa bits should depend on the data. It offers an alternative explanation, based on range and precision alone, which we should keep in mind alongside the EoS view.

\newpage
\section{Implications for machine learning theory and practice}

\paragraph{1. Stability theory must model the arithmetic, not only the step size.}
EoS theory \citep{Cohen2021,Damian2023,Cohen2022} assumes exact gradients, so its threshold is $2/\eta$ (or $38/\eta$ preconditioned) for every format. A format can only matter through the error it injects, and what matters is the error's component along $q_1$ and its bias, not its average size. Su et al.'s multiplicative model lowers the threshold to $2/((1+\|\zeta\|)\eta)$. With shared block scales the error is correlated within each block, so its projection on $q_1$ can be much larger than isotropic-noise models predict. This yields a falsifiable prediction. If EoS explains the gap, E2M5 should equilibrate at a measurably lower sharpness (or preconditioned sharpness) than E3M4 at the same $\eta$. If both settle at the same threshold, curvature is not the channel, and the gap must come from range loss, mismatch, or a slower-converging but equally stable trajectory. Measuring this requires sharpness computed in high precision (Lanczos with Hessian--vector products in FP32/FP64), since Damian et al.\ note that the dynamics are sensitive to small perturbations and switch to 64-bit gradients to avoid propagating floating-point error.

\paragraph{2. Explaining the E3M4 advantage would turn format choice from a sweep into a prediction.}
Today the ranking of scale formats is purely empirical: \citet{Chmiel2025} find E3M4 ahead of E2M5 without explaining why, and \citet{Kuzmin2022} show that the ranking changes with the bias and the tensor distribution. Our project asks which mechanism produces the gap, and each answer has a different consequence. If the gap comes from range, the ranking should follow from tensor statistics, such as the spread of block maxima that a scale must encode, and could be predicted before training. If it comes from the edge of stability, the ranking also depends on the learning rate, and should shrink or vanish when $\eta\lambda_{\max}$ stays well below threshold. For theory, either answer replaces a sweep with a principle that can be tested on new formats. For practice, it tells hardware designers which bits to spend: Chmiel et al.'s results already challenge the MX standard's E8M0 scale, with no mantissa bits, and a mechanism would say whether E3M4 is the right replacement for every model or only for some training regimes.

\paragraph{3. Optimizer-state precision is a stability parameter, not only a memory saving.}
Under Adam, the edge of stability is set by the preconditioned sharpness $\lambda_{\max}(P^{-1}H)$ \citep{Cohen2022}, so any error in $P=\mathrm{diag}(\sqrt{\nu})+\epsilon I$ moves the effective threshold. Memory-efficient training stores $\nu$ in 8 bits \citep{Dettmers2022b}, yet $\nu$ spans many orders of magnitude across parameters. A format with too little range flushes the small entries of each block to zero, so $P$ falls to $\epsilon I$ in those coordinates, their steps grow by orders of magnitude, and $\lambda_{\max}(P^{-1}H)$ can cross $38/\eta$ at a learning rate that is otherwise stable. This gives a second, concrete route by which E3M4 could beat E2M5: through the optimizer state rather than the gradients. It is testable by quantizing only $\nu$ in each format and tracking $\lambda_{\max}(P^{-1}H)$ against $38/\eta$. In practice, $\nu$ should be kept in a range-rich format, or its underflow rate monitored.

\paragraph{4. Practical consequences for low-precision training and RL.}
Because the stable step size can depend on the format, comparing formats at one shared learning rate confounds format with $\eta\lambda$. Learning rates and warmup should be tuned per format, or at least swept. Following \citet{Qi2025}, rollout and training engines should share arithmetic, or their discrepancy should be monitored, whenever probabilities are compared across engines. Finally, monitoring sharpness online is affordable with a few Hessian--vector products every few hundred steps, and could serve as an early warning before the late divergences reported by \citet{Fishman2025}.

\clearpage
\bibliographystyle{plainnat}
{\footnotesize
\begin{thebibliography}{29}
\setlength{\itemsep}{0pt}

\bibitem[Ahn et~al.(2022)]{Ahn2022}
K.~Ahn, J.~Zhang, and S.~Sra.
\newblock Understanding the unstable convergence of gradient descent.
\newblock In \emph{ICML}, 2022. arXiv:2204.01050.

\bibitem[Arora et~al.(2022)]{Arora2022}
S.~Arora, Z.~Li, and A.~Panigrahi.
\newblock Understanding gradient descent on edge of stability in deep learning.
\newblock In \emph{ICML}, 2022. arXiv:2205.09745.

\bibitem[Chmiel et~al.(2025)]{Chmiel2025}
B.~Chmiel, M.~Fishman, R.~Banner, and D.~Soudry.
\newblock {FP4} all the way: Fully quantized training of large language models.
\newblock In \emph{NeurIPS}, 2025. arXiv:2505.19115.

\bibitem[Cohen et~al.(2021)]{Cohen2021}
J.~M. Cohen, S.~Kaur, Y.~Li, J.~Z. Kolter, and A.~Talwalkar.
\newblock Gradient descent on neural networks typically occurs at the edge of stability.
\newblock In \emph{ICLR}, 2021. arXiv:2103.00065.

\bibitem[Cohen et~al.(2022)]{Cohen2022}
J.~M. Cohen, B.~Ghorbani, S.~Krishnan, N.~Agarwal, S.~Medapati, M.~Badura, D.~Suo, D.~Cardoze, Z.~Nado, G.~E. Dahl, and J.~Gilmer.
\newblock Adaptive gradient methods at the edge of stability.
\newblock arXiv:2207.14484, 2022.

\bibitem[Damian et~al.(2023)]{Damian2023}
A.~Damian, E.~Nichani, and J.~D. Lee.
\newblock Self-stabilization: The implicit bias of gradient descent at the edge of stability.
\newblock In \emph{ICLR}, 2023. arXiv:2209.15594.

\bibitem[Dettmers et~al.(2022a)]{Dettmers2022}
T.~Dettmers, M.~Lewis, Y.~Belkada, and L.~Zettlemoyer.
\newblock {LLM.int8()}: 8-bit matrix multiplication for transformers at scale.
\newblock In \emph{NeurIPS}, 2022. arXiv:2208.07339.

\bibitem[Dettmers et~al.(2022b)]{Dettmers2022b}
T.~Dettmers, M.~Lewis, S.~Shleifer, and L.~Zettlemoyer.
\newblock 8-bit optimizers via block-wise quantization.
\newblock In \emph{ICLR}, 2022. arXiv:2110.02861.

\bibitem[Fishman et~al.(2025)]{Fishman2025}
M.~Fishman, B.~Chmiel, R.~Banner, and D.~Soudry.
\newblock Scaling {FP8} training to trillion-token {LLMs}.
\newblock In \emph{ICLR}, 2025. arXiv:2409.12517.

\bibitem[Gupta et~al.(2015)]{Gupta2015}
S.~Gupta, A.~Agrawal, K.~Gopalakrishnan, and P.~Narayanan.
\newblock Deep learning with limited numerical precision.
\newblock In \emph{ICML}, 2015.

\bibitem[He and Thinking Machines Lab(2025)]{He2025}
H.~He and Thinking Machines Lab.
\newblock Defeating nondeterminism in {LLM} inference.
\newblock \emph{Thinking Machines Lab: Connectionism}, 2025. \url{https://thinkingmachines.ai/blog/defeating-nondeterminism-in-llm-inference/}.

\bibitem[Jastrz{\k{e}}bski et~al.(2020)]{Jastrzebski2020}
S.~Jastrz{\k{e}}bski, M.~Szymczak, S.~Fort, D.~Arpit, J.~Tabor, K.~Cho, and K.~Geras.
\newblock The break-even point on optimization trajectories of deep neural networks.
\newblock In \emph{ICLR}, 2020. arXiv:2002.09572.

\bibitem[Kreisler et~al.(2023)]{Kreisler2023}
I.~Kreisler, M.~S. Nacson, D.~Soudry, and Y.~Carmon.
\newblock Gradient descent monotonically decreases the sharpness of gradient flow solutions in scalar networks and beyond.
\newblock In \emph{ICML}, 2023. arXiv:2305.13064.

\bibitem[Kuzmin et~al.(2022)]{Kuzmin2022}
A.~Kuzmin, M.~van Baalen, Y.~Ren, M.~Nagel, J.~Peters, and T.~Blankevoort.
\newblock {FP8} quantization: The power of the exponent.
\newblock In \emph{NeurIPS}, 2022. arXiv:2208.09225.

\bibitem[Micikevicius et~al.(2022)]{Micikevicius2022}
P.~Micikevicius, D.~Stosic, N.~Burgess, et~al.
\newblock {FP8} formats for deep learning.
\newblock arXiv:2209.05433, 2022.

\bibitem[Noune et~al.(2022)]{Noune2022}
B.~Noune, P.~Jones, D.~Justus, D.~Masters, and C.~Luschi.
\newblock 8-bit numerical formats for deep neural networks.
\newblock arXiv:2206.02915, 2022.

\bibitem[Qi et~al.(2025)]{Qi2025}
P.~Qi, Z.~Liu, X.~Zhou, T.~Pang, C.~Du, W.~S. Lee, and M.~Lin.
\newblock Defeating the training-inference mismatch via {FP16}.
\newblock arXiv:2510.26788, 2025.

\bibitem[Darvish~Rouhani et~al.(2023)]{Rouhani2023}
B.~Darvish~Rouhani, R.~Zhao, A.~More, et~al.
\newblock Microscaling data formats for deep learning.
\newblock arXiv:2310.10537, 2023.

\bibitem[Su et~al.(2025)]{Su2025}
H.~Su, M.~Kwun, S.~Gil, S.~Kakade, and N.~Anand.
\newblock Characterization and mitigation of training instabilities in microscaling formats.
\newblock arXiv:2506.20752, 2025.

\bibitem[Sun et~al.(2019)]{Sun2019}
X.~Sun, J.~Choi, C.-Y. Chen, et~al.
\newblock Hybrid 8-bit floating point ({HFP8}) training and inference for deep neural networks.
\newblock In \emph{NeurIPS}, 2019.

\bibitem[van Baalen et~al.(2023)]{vanBaalen2023}
M.~van Baalen, A.~Kuzmin, et~al.
\newblock {FP8} versus {INT8} for efficient deep learning inference.
\newblock arXiv:2303.17951, 2023.

\bibitem[Wortsman et~al.(2023)]{Wortsman2023}
M.~Wortsman, T.~Dettmers, L.~Zettlemoyer, A.~Morcos, A.~Farhadi, and L.~Schmidt.
\newblock Stable and low-precision training for large-scale vision-language models.
\newblock In \emph{NeurIPS}, 2023. arXiv:2304.13013.

\bibitem[Wu et~al.(2018)]{Wu2018}
L.~Wu, C.~Ma, and W.~E.
\newblock How {SGD} selects the global minima in over-parameterized learning: A dynamical stability perspective.
\newblock In \emph{NeurIPS}, 2018.

\bibitem[Yao et~al.(2025)]{Yao2025}
F.~Yao, L.~Liu, D.~Zhang, C.~Dong, J.~Shang, and J.~Gao.
\newblock Your efficient {RL} framework secretly brings you off-policy {RL} training.
\newblock Blog post, 2025. \url{https://fengyao.notion.site/off-policy-rl}.

\end{thebibliography}}
\end{document}
